\begin{onlytopic}
	\chapter{小数计算}
\end{onlytopic}
%%%%%%练习%%%%%%%%%%%%%%%%%%%%%%%%%%%%%%

\begin{zuoye}
	\begin{problems}
		计算 ：		
		$12.7+5.63=$ $\_\_\_\_$\hfil $6-3.125=$ $\_\_\_\_$ 
		
		 $0.35 \times 0.4=$ $\_\_\_\_$ ．\hfil  $9.999 \div 3.3=$ $\_\_\_\_$ ．
	\end{problems} 
	
	\begin{lxsol}
		\vspace{\solthreeproblems}
	\end{lxsol}
\end{zuoye}

\begin{zuoye}
	\begin{problems}
		计算：$3.71-2.74+4.7+5.29-0.26+6.3$
	\end{problems} 
	
	\begin{lxSsol}
		\vspace{\solthreeproblems}
	\end{lxSsol}
\end{zuoye}

\begin{zuoye}
	\begin{problems}
		计算 ： $0.2+0.4+0.6+0.8+0.12+0.14+0.16+0.18=$ $\_\_\_\_$ ．
	\end{problems} 
	
	\begin{lxSsol}
		\vspace{\solthreeproblems}
	\end{lxSsol}
\end{zuoye}

\begin{zuoye}
	\begin{problems}
		计算 (1)$4 \div 11 \times 3.3 \times 2.5=$ $\_\_\_\_$\hfil  (2) $0.125 \times 7.2 \div 0.9=$ $\_\_\_\_$
		
		(3)$11 \div 0.125 \div 8=$ $\_\_\_\_$ ．
	\end{problems} 
	
	\begin{lxMsol}
		\vspace{\solthreeproblems}
	\end{lxMsol}
\end{zuoye}


\begin{zuoye}
	\begin{problems}
		计算  (1)$8.3 \times 12.6-2.6 \times 8.3$ ．\hfil  (2)$ 67 \times 8.1+67 \times 10.1+67 \times 12.1-67 \times 0.3$	
	\end{problems} 
	
	\begin{lxMsol}
		\vspace{\solthreeproblems}
	\end{lxMsol}
\end{zuoye}




\begin{zuoye}
	\begin{problems}
		计算： $2 \times 0.3 \times 5 \times 7 \times 1.1 \times 1.3 \times 1.7 \times 1.9 \div 3.8 \div 0.51 \div 6.5 \div 7.7$
	\end{problems} 
	
	\begin{lxMsol}
		\vspace{\solthreeproblems}
	\end{lxMsol}
\end{zuoye}
%%%%%%练习%%%%%%%%%%%%%%%%%%%%%%%%%%%%%%


%%%%%%练习%%%%%%%%%%%%%%%%%%%%%%%%%%%%%%

\begin{zuoye}
	计算: $20.09 \times 31.5+2.009 \times 317+200.9 \times 3.68$ \hfil $0.888 \times 125 \times 73+999 \times 3$
	
	\begin{lxMsol}
		\vspace{\solfourproblems}
	\end{lxMsol}
\end{zuoye}

%%%%%%%%%%%


%%%%%%%%%%%
\begin{zuoye}
	\begin{problems}
	已知等式 $105.5+[(40+\square \div 2) \times 0.5-1.5] \div(53.6 \div 26.8 \times 0.125)=185.5$ ，那么$\square =$ \tk ．
	\end{problems} 
	
	\begin{lxMsol}
		\vspace{\solthreeproblems}
	\end{lxMsol}
\end{zuoye}
%%%%%%练习%%%%%%%%%%%%%%%%%%%%%%%%%%%%%%

\begin{zuoye}
	两位带小数相乘, 积四舍五入后是 $39.1$, 这两个数都是一位小数, 且个位 上的数字都是 6, 那么乘积四舍五入前是几?	
	
	\begin{lxSsol}
		\vspace{\solfourproblems}
	\end{lxSsol}
\end{zuoye}

%%%%%%%%%%%


%%%%%%练习%%%%%%%%%%%%%%%%%%%%%%%%%%%%%%
%%
%\begin{zuoye}
%	从 $19,20,21, \ldots, 92,93,94$ 这 76 个数中, 选取两个不同的数, 使其和为偶 数的选法总数是多少?
%	
%	\begin{lxMsol}
%		\vspace{\solfourproblems}
%	\end{lxMsol}
%\end{zuoye}

%%%%%%%%%%%

%
%
%
%%%%%%%练习%%%%%%%%%%%%%%%%%%%%%%%%%%%%%%
%
%\begin{zuoye}
%	有 5 人参加的学雷锋小队上街宣传交通规则, 他们站成一排, 其中 2 名队长不排在一起, 一共有多少种排法.	
%	
%	
%	\begin{lxMsol}
%		\vspace{\solfourproblems}
%	\end{lxMsol}
%\end{zuoye}
%
%%%%%%%%%%%%
%
%
%%%%%%%练习%%%%%%%%%%%%%%%%%%%%%%%%%%%%%%
%
%\begin{zuoye}
%	如图, 把 $A , B , C , D , E$ 这五部分用四种不同的颜色着色, 且相邻的部 分不能使用同一种颜色, 不相邻的部分可以使用同一种颜色.那么, 这幅图一共 有多少种不同的着色方法?	
%	
%	
%	\begin{figure}[H]
%		\begin{flushright}
%			\begin{tikzpicture}[scale=0.6]
%				\draw (0,0) rectangle(3,4);
%				\draw (1,0) --(1,3) node at (0.5,0.5){$D$};
%				\draw (0,3) --(3,3) node at (2,0.5){$E$};
%				\draw (1,1) --(3,1) node at (2,2){$C$};
%				\draw (1,2) --(0,2) node at (0.5,2.5){$B$};
%				\node at (1.5,3.5) {$A$};
%			\end{tikzpicture}
%		\end{flushright}
%		
%	\end{figure}
%	\begin{lxsol}
%		\vspace{\solfiveproblems}
%	\end{lxsol}
%\end{zuoye}

%%%%%%%
%
%小测
%\begin{zuoye}
%	计算:  $3.49+4.47+3.51+2.38+4.53+4.62$	
%	
%	\begin{lxMsol}
%		\vspace{\solfourproblems}
%	\end{lxMsol}
%\end{zuoye}
%$1.242 \times 6.7+0.758 \times 6.7$ 